\section{Algorithm}
\label{sec:alg}

This section describes our parallel congruence closure procedure in
detail. The algorithm rests on three components: a lock-free concurrent
union-find that allows arbitrarily many threads to issue
\textsc{find}/\textsc{union} concurrently
(Section~\ref{subsec:uf}); a parallel semisort that lets each round
group congruent e-node candidates by signature without a sequential
hash-table bottleneck (Section~\ref{subsec:semisort}); and a
bulk-synchronous, rounds-based outer loop that drives both to a
congruence-closed fixed point (Section~\ref{subsec:rounds}). All three
build on \textsc{ParlayLib}~\cite{parlaylib}, a shared-memory
work-stealing runtime that supplies the sequence types and parallel
primitives (\texttt{parallel\_for}, \texttt{filter},
\texttt{flat\_map}, \texttt{group\_by\_key}) that we use throughout.

\subsection{Concurrent Union-Find}
\label{subsec:uf}

We use the lock-free concurrent disjoint-set-union of Jayanti and
Tarjan~\cite{jayanti16}. A union-find of capacity $n$ is represented as
an array of $n$ atomic 32-bit words. Each word encodes either a
\emph{rank} (when the node is a root, indicated by the high bit) or a
\emph{parent pointer} into the same array (when the node is not a
root). Both \textsc{find} and \textsc{union} are non-blocking: they
read the current state, decide what update they would like to make,
and install it via a compare-and-swap (CAS); on CAS failure they retry
from a freshly-read state.

\textsc{find}$(u)$ recursively walks parent pointers to a root, and on
its way back attempts to install path compression by CAS-ing the
walked node's parent slot directly to the root. Because the CAS only
succeeds when the slot still points to the old parent, lost races are
harmless: another thread has already installed an at-least-as-good
shortcut. \textsc{union}$(u, v)$ first calls \textsc{find} on both
endpoints, then attempts to CAS the smaller-rank root's slot from a
rank value to a parent pointer to the other root, retrying from the
top of the routine on failure. Equal ranks tie-break by node id and
attempt a best-effort rank increment on the winning root with a second
CAS.

This UF supports the throughput regime that our outer loop demands:
each round of congruence closure issues thousands to millions of
\textsc{find} and \textsc{union} calls in parallel, and every operation
linearizes regardless of contention.

\subsection{Parallel Semisort}
\label{subsec:semisort}

At each round, we want to group congruent terms based on newly discovered equalities. To this end, we use the \emph{semisort} of Gu, Shun, Sun, and Blelloch~\cite{semisort}. Given a type $K$, and a sequence $S$ of $K$, an equality comparator $\mathit{eq} : K \times K \to \mathbb{B}$, and a hash function $h : K \to \mathbb{N}$, semisort will group equal keys in this sequence adjacently. Concretely, we use semisort to group congruence classes at each round.

\subsection{Rounds-Based Congruence Closure}
\label{subsec:rounds}

At each round, we maintain a set of terms whose congruence classes may need to be updated. Intially, this comprises all terms appearing in base equalities. We then group congruent terms based on the equalities currently discovered, and merge the congruence classes of each group in the union-find. Grouping is accomplished via semisort---the \texttt{group\_by} operation in \texttt{parlay}---and merging congruence classes is accomplished through taking the union of the corresponding equality classes in the union-find.

We then propagate terms whose congruence classes may have changed to the next round. These are exactly the \emph{predecessors} of nodes that whose congruence classes have changed that round. Specifically, if the congruence class of term $e$ has changed, then the congruence class of $f(e)$ may have changed as well. This requires maintaining the predecessors of every congruence class; specifically a mapping \textit{pred} between class identifiers and identifiers of predecessors of that class. This mapping must be updated every round; if two congruence classes are merged, then their predecessors must be merged as well. These rounds continue until there is no longer any work to do. Pseudocode depicting this algorithm is shown in Figure~\ref{fig:alg:bsp}.

The algorithm is decomposed into three different functions. The function \textsc{Congruent} checks whether two terms are congruent based on the equivalence classes in the union find. Two terms are congruent if they are both the same function symbol and each of their subterms are pointwise congruent.

The \textsc{Merge} function is used to merge the equivalence classes of different terms. Specifically, given a sequence $S$ of terms, it will union the classes of every term in $S$. We take a simple divide-and-conquer approach; we merge the left and right halves (in parallel), and then \textsc{Union} the representatives of the two classes.

\begin{figure}
  \centering
\begin{algorithmic}[1]
    \Function{Congruent}{X, Y}
      \If{$X = f(M_1,\ldots,M_n)$ and $Y = f(N_1,\ldots,N_n)$}
        \State \Return $\Call{Find}{M_i} = \Call{Find}{N_i}$ for all $i = 1,\ldots,n$
      \Else 
        \State \Return \textbf{false}
      \EndIf
    \EndFunction

    \Function{MergeCongruenceClass}{\textit{Class}}
      \If{$\mathit{Class} = \{e\}$}
        \State \Return{$e$} \Comment{Singleton class, return representative}
      \EndIf
      \State $n \gets |\textit{Class}|$
      \State $L \gets \mathit{Class}[0,\ldots,\frac{n}{2}]$
      \State $R \gets \mathit{Class}[\frac{n}{2},\ldots,n]$
      \State $\rhd$ Merge left and right halves, get their representatives
      \State $l, r \gets \Call{MergeCongruenceClass}{L} \parallel \Call{MergeCongruenceClass}{R}$
      \State \Return $\Call{Union}{l, r}$ \Comment{Union left and right reps, return root}
    \EndFunction

    \Function{ParallelClose}{equalities = $\{e_1 = e_1', \ldots, e_w = e_w'\}$}
        \State $\mathit{Frontier} \gets \text{every term appearing in } \mathit{equalities}$
        \While{$\mathit{Frontier}$ is not empty}
          \State $\mathit{groups} \gets \text{congruence classes quotiented by } \textsc{Congurent}$.
          \ParFor{$\mathit{group} \in \mathit{groups}$}
            \State $\rhd$ Merge every term within congruence class
            \State $\mathit{rep} \gets \Call{Merge}{\mathit{group}}$
            \State $\rhd$ Construct new parent lists
            \State $\mathit{groups}[\mathit{rep}] \gets \Call{ConcatMap}{\lambda \mathit{group}. \; \mathit{Parents}[\mathit{group}]}$
          \EndParFor
          \State $\mathit{Frontier} \gets \Call{Flatten}{\mathit{groups}}$
        \EndWhile
    \EndFunction
\end{algorithmic}

\caption{Round-based congruence closure algorithm}
\label{fig:alg:bsp}
\end{figure}

% \begin{algorithmic}[1]
%     \Function{Merge}{M, N}
%         \If{$\Call{Find}{M} = \Call{Find}{N}$}
%             \State \Return
%         \EndIf
%         \State \Call{Union}{M,N}
%         \State Let $P_M$ be the predecessors associated with the congruence class of $M$
%         \State Let $P_N$ be the predecessors associated with the congruence class of $N$
%         \For{$(X, Y) \in P_M \times P_N$}
%           \If{$\Call{Find}{X} \neq \Call{Find}{Y}$ and $\Call{Congruent}{X, Y}$}
%             \State $\Call{Merge}{X, Y}$
%           \EndIf
%         \EndFor
%     \EndFunction
%     \Function{Congruent}{X, Y}
%       \If{$X = f(M_1,\ldots,M_n)$ and $Y = f(N_1,\ldots,N_n)$}
%         \State \Return $\Call{Find}{M_i} = \Call{Find}{N_i}$ for all $i = 1,\ldots,n$
%       \Else 
%         \State \Return \textbf{false}
%       \EndIf
%     \EndFunction
% \end{algorithmic}

% \begin{algorithm}
% \caption{Parallel congruence closure (one round)}
% \label{alg:rounds}
% \textbf{Input:} 

% e-graph $G$ with concurrent UF $U$, parent-index
% $P[\,\cdot\,]$, per-class \emph{changed} flags $C[\,\cdot\,]$, after a
% batch of initial unions has marked each affected class. \\
% \textbf{Output:} A new round's worth of merges; loop until no merges
% fire.
% \begin{enumerate}
%   \item \emph{Frontier.} In parallel over class ids, \texttt{filter} $C$
%         to collect $F = \{c : C[c]\}$ and clear the corresponding flags.
%   \item \emph{Gather.} \texttt{flat\_map} $F$ through $P$ to obtain the
%         list of candidate parent e-nodes whose signature might have
%         shifted.
%   \item \emph{Canonicalize.} \texttt{map} each candidate $n$ to a
%         triple $(h, \mathit{idx}, r)$ where $h$ is a 64-bit hash of
%         $(\mathrm{op}(n), U.\textsc{find}(c_1), \dots,
%                        U.\textsc{find}(c_k))$,
%         $\mathit{idx}$ is the candidate's position in the e-node array,
%         and $r = U.\textsc{find}(\mathrm{class}(n))$.
%   \item \emph{Group.} Hash $h$ into one of $B$ buckets and apply the
%         parallel semisort of Section~\ref{subsec:semisort}.
%   \item \emph{Detect merges.} In parallel over buckets, sort within
%         each bucket by $(h, \text{lexicographic canonicalized
%         children})$ and scan adjacent pairs: every pair with equal
%         signature whose representatives differ contributes a
%         $(r_a, r_b)$ pair to a new \emph{merge list}.
%   \item \emph{Apply.} In parallel over the merge list, call
%         $U.\textsc{union}(r_a, r_b)$ and set $C[r_a] = C[r_b] =
%         \texttt{true}$.
% \end{enumerate}
% \end{algorithm}
